\documentclass[11pt]{article}

\usepackage[margin=1in]{geometry}
\usepackage{amsmath,amssymb}

\title{Preprocessing Pipeline for Monthly Return Prediction}
\date{}

\begin{document}
\maketitle

\section{Notation}

Let $i \in \{1,\dots,N_c\}$ index companies and let $t$ index months. Each row in the
raw file corresponds to a company--month pair $(i,t)$. We denote the target by
$y_{i,t}$ (the column \texttt{monthly\_gross\_return}), a gross return where values
$> 1$ indicate gains and values $< 1$ indicate losses. For interpretation one may
define the net return $r_{i,t} = y_{i,t} - 1$, but the pipeline operates on $y_{i,t}$.
The raw feature vector for a row is $\mathbf{x}_{i,t} \in \mathbb{R}^d$. Define the
design matrix $X \in \mathbb{R}^{N \times d}$ where row index $r = r(i,t)$ maps to a
specific pair $(i,t)$, so $X_{r,:} = \mathbf{x}_{i,t}$. Missingness indicators are
binary features $m_{i,t,j} = \mathbb{1}[x_{i,t,j}\ \text{is missing}]$.

\section{Raw Data Schema}

The raw data are a panel of company--month observations. Key fields are:
\begin{itemize}
	\item \textbf{Identifiers:} \texttt{co\_code} (company id) and \texttt{Month}
	(YYYY-MM timestamp).
	\item \textbf{Target:} \texttt{monthly\_gross\_return}.
	\item \textbf{Market variables:} \texttt{mktcap}, \texttt{price},
	\texttt{lag\_mv}, \texttt{BM\_sep}.
	\item \textbf{Accounting variables:} \texttt{equity\_bv\_on\_stkdate},
	\texttt{fy\_book\_value}, \texttt{sa\_net\_worth}, \texttt{sa\_total\_assets},
	\texttt{sa\_sales}, \texttt{sa\_cost\_of\_goods\_sold},
	\texttt{sa\_total\_interest\_exp}, \texttt{sa\_selling\_distribution\_exp},
	\texttt{lagged\_book\_equity}, \texttt{lagged\_assets}, \texttt{lagged2\_assets}.
	\item \textbf{Factor-style signals:} \texttt{OpProf}, \texttt{Inv},
	\texttt{Momentum}.
	\item \textbf{Categorical labels:} \texttt{Size\_Label} (S/B), \texttt{BM\_Label}
	(G/N/V), \texttt{OpProf\_Label} (W/N/R), \texttt{Inv\_Label} (A/N/C),
	\texttt{Mom\_Label} (L/N/W).
	\item \textbf{Calendar support:} \texttt{Year}, \texttt{Corrected\_Year},
	\texttt{Corrected\_Month}, \texttt{Month\_annual}.
\end{itemize}
Any remaining columns are treated as numeric features.

\section{Pipeline Overview}

The preprocessing mapping is
\[
\texttt{factor\_data.csv}
\xrightarrow{\text{type coercion + cleaning}}
\xrightarrow{\text{lag/rolling features}}
\xrightarrow{\text{missing indicators}}
\xrightarrow{\text{split assignment}}
\{\texttt{train},\texttt{validation},\texttt{test}\}.
\]

\section{Type Coercion and Cleaning}

\begin{itemize}
	\item \textbf{Column normalization:} column names are stripped of whitespace.
	\item \textbf{Date parsing:} \texttt{Month} is parsed to a datetime. Rows with
	invalid \texttt{Month} are dropped.
	\item \textbf{Identifiers:} \texttt{co\_code} is coerced to nullable integer.
	\item \textbf{Target:} \texttt{monthly\_gross\_return} is coerced to float.
	\item \textbf{Categorical labels:} categorical columns are coerced to strings,
	stripped, and empty tokens (\texttt{empty string}, \texttt{nan}, \texttt{<NA>})
	are converted to missing values.
	\item \textbf{Numeric features:} all remaining non-categorical, non-identifier
	columns are coerced to floats; non-numeric values become missing.
\end{itemize}

Raw return columns \texttt{lag\_ret} and \texttt{log\_ret} are removed to avoid
leakage. All engineered return history is derived from the target $y_{i,t}$.

\section{Leakage-Safe Feature Engineering}

All time features are computed within each company after sorting by
$(\texttt{co\_code}, \texttt{Month})$.

\subsection{Lag Features}

For lag steps $k \in \{1,3,6\}$,
\[
\text{ret\_lag}_{k}(i,t) = y_{i,t-k}.
\]
These lags use only past values and therefore respect causal ordering.

\subsection{Rolling Statistics}

Let $w \in \{3,6\}$ denote the rolling window length and define the available
past returns as $\{y_{i,t-1}, y_{i,t-2}, \dots\}$. For each window, the pipeline
computes
\[
\mu^{(w)}_{i,t} = \frac{1}{m} \sum_{s=1}^{m} y_{i,t-s},
\quad m = \min\{w,\ \#\text{available past returns}\},
\]
and
\[
\sigma^{(w)}_{i,t} = \sqrt{\frac{1}{m-1} \sum_{s=1}^{m} (y_{i,t-s}-\mu^{(w)}_{i,t})^2},
\quad m \ge 2.
\]
The script also records the count of available values in the 3-month window,
\[
\text{ret\_rolling\_count\_3m}(i,t) = \#\{s \in \{1,2,3\} : y_{i,t-s}\ \text{observed}\}.
\]

\subsection{Month Index}

The monotone calendar index is
\[
\text{month\_index}_{i,t} = 12 \cdot \text{year}(t) + \text{month}(t).
\]

\section{Missingness Indicators}

For every column except \texttt{co\_code} and \texttt{Month}, a missingness indicator
is added:
\[
m_{i,t,j} = \mathbb{1}[x_{i,t,j}\ \text{is missing}].
\]
This keeps the original missing values while making their presence explicit to
downstream models.

\section{Split Assignment}

Two split modes are supported.

\subsection{Time-Based Split}

Let $t^*$ be the split date for each row. If \texttt{split\_date\_source=corrected}
and corrected fields exist, then $t^*$ is formed from
\texttt{Corrected\_Year} and \texttt{Corrected\_Month}; otherwise $t^* = \texttt{Month}$.
Given cutoffs $\tau_{\text{train}}$ and $\tau_{\text{val}}$, define
\[
\text{split}(i,t) =
\begin{cases}
\text{train}, & t^* \le \tau_{\text{train}},\\
\text{validation}, & \tau_{\text{train}} < t^* \le \tau_{\text{val}},\\
\text{test}, & t^* > \tau_{\text{val}}.
\end{cases}
\]

\subsection{Company-Based Split}

Companies are assigned to splits in aggregate. Let $H_i$ be the number of
distinct months for company $i$. The validation set is drawn first from the
eligible set $\{i : H_i \ge H_{\min}\}$; if insufficient, the remaining slots are
filled with shorter-history companies. Remaining companies are split into train
and test to match the requested shares. A fixed random seed ensures
reproducibility and each company appears in exactly one split.

\section{Leakage-Safe Test Outputs}

For the test split, the target and raw return columns are masked to avoid
inference leakage:
\[
y_{i,t} \leftarrow \text{NaN} \quad \text{for test rows}.
\]
The original values are written to \texttt{test\_targets.csv} for evaluation. Any
missing indicators for masked columns are updated to reflect the new missing
status.

\section{Holdout Company (\texttt{co\_code = 194})}

In addition to the train/validation/test splits, company \texttt{194} is removed
from the initial dataset and treated as a standalone holdout. This produces a
clean, out-of-sample company-level evaluation set that never appears in model
training or split assignment. The holdout rows follow the same preprocessing
steps (type coercion, leakage-safe feature engineering, and missingness
indicators) before inference and reporting.

\section{Outputs}

The script writes:
\begin{itemize}
	\item \texttt{processed\_full.csv}: full dataset with engineered features and
	\texttt{split\_group} labels.
	\item \texttt{train.csv}, \texttt{validation.csv}, \texttt{test.csv}: split
	subsets, with the test target masked.
	\item \texttt{test\_targets.csv}: true target (and any raw return columns if
	present) for the test rows.
	\item \texttt{metadata.json}: row counts, split counts, categorical columns,
	and a feature column list (excluding identifiers, dates, target, and
	missing-indicator columns).
\end{itemize}

The resulting tables are ready for training and inference: CatBoost consumes the
tabular features directly, while the CNN stage builds short temporal sequences
from the same leakage-safe history.

\end{document}
